\section{多尺度探索与跨任务泛化}
\subsection{多阶段探索策略}
Slide 07 展示 multi-stage exploration framework：将 exploration policy 分成 coarse-grained 与 fine-grained 两层，分别控制不同尺度的目标与动作序列。这样的层次结构既能发现 macro-level goal cluster，又能自动收敛到 micro-level 执行步骤。

\begin{figure}[H]
\centering
\includegraphics[width=0.85\textwidth]{slides-images/slide-07.jpg}
\caption{Slide 7 说明 multi-stage exploration 如何协同扩展 goal space。}
\end{figure}

\begin{knowledgebox}{Multi-Stage Exploration}
\begin{itemize}
\item Coarse policy 采样未见过的 goal cluster；\\
\item Fine policy 负责该 cluster 内的精细动作；\\
\item 定期更新 cluster centroids based on success rate 以维持 coverage。
\end{itemize}
\end{knowledgebox}

\subsection{跨任务泛化指标}
Slide 08 提出 transfer matrix 检验 policy 在不同 goal 组合上的 capability，关键指标包括 success coverage、policy entropy 与 transfer ratio。讲者建议在实验报告中同时展示 base goals 与 transfer goals 的 performance，避免过拟合单一场景。

\begin{tabularx}{\textwidth}{lX}
\toprule
指标 & 说明 \\
\midrule
Success coverage & 在 stored goals 中达到成功的比例，低于 70\% 表明 coverage gap； \\
Policy entropy & 监控 policy 在不同 goals 上的多样性，过低说明 degenerate； \\
Transfer ratio & 一个 goal 的 reward 在另一 goal 上带来的 improvement，衡量泛化； \\
\bottomrule
\end{tabularx}

\subsection{本章小结}
多尺度探索与泛化指标帮助团队识别在哪些 goal 上要打破窄瓶颈，并提供衡量 transfer 成效的工具。
\newpage

